\section{Results}
\label{sec:results}

\subsection{Rebalancing and unconstrained augmentation are insufficient}
\label{sec:dacpf-results}

The project began from the conventional explanation for poor rare-outcome prediction: that the minority class simply supplies too few examples. We first examined ordinary rebalancing, CTGAN minority augmentation, and increasingly targeted candidate filters. The common result across these attempts was not that augmentation never helped, but that additional minority records did not produce a sufficiently large and consistent gain to explain the problem by class frequency alone. Some of these early notebooks were single-split or diagnostic studies and are treated as hypothesis-generating rather than confirmatory evidence.

The leakage-controlled \dacpf{} rerun provides the strongest five-fold evidence from this stage (Table~\ref{tab:dacpf}). Standard CTGAN augmentation improves mean Macro-F1 relative to real-only training for both classifiers available in the local environment. Adding an age filter, or a combined age-and-disease filter, does not improve on standard CTGAN. Disease-anchor filtering performs best, adding 0.0508 over standard CTGAN for CatBoost but only 0.0084 for XGBoost. Domain targeting can therefore help, but its effect depends on both the anchor and the classifier, and adding more constraints is not monotonically beneficial.

\begin{table}[htbp]
\centering
\caption{Five-fold COVID-19 augmentation Macro-F1. These are Real+Synthetic$\rightarrow$Real results and are not directly comparable with the standalone-synthetic \tstr{} results in Section~\ref{sec:fixed-lambda}.}
\label{tab:dacpf}
\begin{tabular}{lrr}
\toprule
Condition & CatBoost & XGBoost\\
\midrule
Real only & 0.4740 & 0.4759\\
Real + standard CTGAN & 0.5262 & 0.5255\\
Real + age-filtered CTGAN & 0.5257 & 0.4983\\
Real + disease-filtered CTGAN & \textbf{0.5771} & \textbf{0.5339}\\
Real + combined \dacpf{} & 0.5256 & 0.5204\\
\bottomrule
\end{tabular}
\end{table}

\begin{figure}[htbp]
\centering
\includegraphics[width=.75\textwidth]{covid_dacpf_summary.pdf}
\caption{COVID-19 \dacpf{} Macro-F1 averaged across CatBoost and XGBoost. The classifier-specific values in Table~\ref{tab:dacpf} show why this average should not be interpreted as a universal anchor effect.}
\label{fig:dacpf}
\end{figure}

This stage ruled out a simple monotonic story in which balancing classes, generating more synthetic cases, or imposing more domain constraints necessarily improves prediction. We then used explainability and neighbourhood analyses to examine diversity, class-wise recall, feature-importance shifts, silhouette, nearest-neighbour label agreement, and hypothesis margin. Because some of these analyses used representative folds or earlier configurations, they establish the diagnostic rationale for what follows rather than a further confirmatory result. Their implication was that weak or unstable class geometry remained after augmentation, which led to the working hypothesis that class overlap, rather than imbalance alone, was limiting synthetic-to-real transfer.

\subsection{Separability as an explicit generation objective}
\label{sec:fixed-lambda}

\sepa{} is the intervention implied by that diagnosis. It does not change the target class counts; instead it asks whether records selected for label-consistent geometry transfer better to held-out real cases. The fixed-$\lambda$ experiment fits the complete pipeline independently in five stratified folds. Table~\ref{tab:fixedutility} reports classifier-averaged fold means and 95\% $t$ intervals. On COVID-19, the balanced setting reaches 0.5492 Macro-F1, but its interval is wide. On Vigitel, both the varied grid and the fixed high-separability setting reach approximately 0.58 and outperform ordinary CTGAN in all ten matched fold--classifier comparisons. The varied-grid result remains exploratory because the configuration was selected using the same folds used for reporting.

\begin{table}[htbp]
\centering
\caption{Five-fold \tstr{} Macro-F1. Means and 95\% intervals are computed from five classifier-averaged fold scores.}
\label{tab:fixedutility}
\begin{tabular}{llrr}
\toprule
Dataset & Training data & Mean & 95\% half-width\\
\midrule
\multirow{4}{*}{COVID} & Real & 0.5069 & 0.0690\\
& Standard CTGAN & 0.4627 & 0.0407\\
& \sepa{} varied grid & 0.5350 & 0.0649\\
& Fixed balanced $(1.0,0.5)$ & \textbf{0.5492} & 0.0959\\
\midrule
\multirow{4}{*}{Vigitel} & Real & 0.4672 & 0.0088\\
& Standard CTGAN & 0.4588 & 0.0107\\
& \sepa{} varied grid & \textbf{0.5839} & 0.0225\\
& Fixed balanced $(1.0,0.5)$ & 0.5834 & 0.0216\\
\bottomrule
\end{tabular}
\end{table}

\begin{figure}[htbp]
\centering
\includegraphics[width=.9\textwidth]{fixed_lambda_fivefold_utility.pdf}
\caption{Five-fold \tstr{} Macro-F1 for the principal fixed-$\lambda$ comparison. Error bars are 95\% $t$ intervals over the five classifier-averaged fold scores.}
\label{fig:fixedutility}
\end{figure}

Table~\ref{tab:lambda} reports the highest mean selection scores from the same run. The repository composite fidelity score varies little across configurations, whereas SI and HM respond strongly to separability pressure. Figure~\ref{fig:lambda-tradeoff} shows the same pattern across the full grid: global fidelity is comparatively stable while separability changes substantially. This is evidence that the selector controls a property not captured by the composite fidelity score; it is not evidence that the selected records are clinically representative.

\begin{table}[htbp]
\centering
\caption{Top executed \sepa{} grid configurations. ``Quality'' is the repository composite fidelity score.}
\label{tab:lambda}
\resizebox{\textwidth}{!}{%
\begin{tabular}{llrrrr}
\toprule
Dataset & $(\lambda_{sep},\lambda_{anchor})$ & Quality & SI & HM & Selection\\
\midrule
COVID & $(1.0,0.0)$ & 0.956 & 0.779 & 0.897 & 2.154\\
COVID & $(1.0,0.25)$ & 0.961 & 0.760 & 0.854 & 2.095\\
COVID & $(1.0,0.5)$ & 0.965 & 0.754 & 0.830 & 2.067\\
Vigitel & $(1.0,0.5)$ & 0.967 & 0.812 & 0.829 & 2.124\\
Vigitel & $(1.0,0.25)$ & 0.967 & 0.811 & 0.822 & 2.116\\
Vigitel & $(1.0,0.0)$ & 0.966 & 0.811 & 0.822 & 2.116\\
\bottomrule
\end{tabular}
}
\end{table}

\begin{figure}[htbp]
\centering
\includegraphics[width=.88\textwidth]{lambda_fidelity_separability.pdf}
\caption{Five-fold mean fidelity--separability trade-off across executed $\lambda$ settings. Colour represents separability weight; marker size represents anchor weight.}
\label{fig:lambda-tradeoff}
\end{figure}

\subsection{Attributing the gain: replicated factorial experiment}
\label{sec:factorial-results}

The fixed-$\lambda$ results raise a sharper question: was the apparent gain produced by separability pressure, anchor pressure, or their interaction? The replicated $2^2$ experiment (Section~\ref{sec:factorial}) answers this component-level question under four predeclared settings. Figure~\ref{fig:foldprofiles} displays all 40 fold responses. COVID-19 profiles cross and show substantial fold-to-fold variation; Vigitel shows a stable step increase when separability is enabled.

\begin{figure}[htbp]
\centering
\includegraphics[width=.95\textwidth]{factorial_fold_profiles.pdf}
\caption{Fold-level \tstr{} Macro-F1 in the replicated $2^2$ \sepa{} experiment. Black lines show condition means.}
\label{fig:foldprofiles}
\end{figure}

\begin{table}[htbp]
\centering
\caption{Mean \tstr{} Macro-F1 by predeclared factorial condition. Values are means of the five exported fold responses.}
\label{tab:conditionmeans}
\begin{tabular}{lrrrr}
\toprule
Dataset & Realism only & Anchor only & Separability only & Both\\
\midrule
COVID & 0.516 & 0.528 & 0.484 & 0.524\\
Vigitel & 0.470 & 0.480 & 0.579 & 0.576\\
\bottomrule
\end{tabular}
\end{table}

The condition means in Table~\ref{tab:conditionmeans} make the contrast explicit. In Vigitel, enabling separability raises the mean from approximately 0.47--0.48 to 0.576--0.579, while adding anchors makes little difference. COVID-19 shows no corresponding ordering: separability-only is the \emph{lowest} condition mean, and the combined condition returns approximately to realism-only performance.

The blocked raw-scale ANOVA in Table~\ref{tab:anova} is the primary component analysis. For COVID-19, none of the treatment effects is significant, and residual error accounts for 81.53\% of treatment-plus-error variation after separating out the fold block. For Vigitel, separability accounts for 96.00\% of that variation and is highly significant ($p=5.03\times10^{-10}$), while anchor and interaction contributions are negligible. The log-scale model yields the same substantive conclusion.

\begin{table}[htbp]
\centering
\caption{Blocked factorial ANOVA: percentage of experimental variation and $p$-values on raw Macro-F1. Percentages exclude fold sum of squares from the denominator.}
\label{tab:anova}
\begin{tabular}{llrr}
\toprule
Dataset & Source & Variation (\%) & $p$\\
\midrule
\multirow{4}{*}{COVID} & Separability & 5.13 & 0.4019\\
& Anchor & 10.21 & 0.2439\\
& Interaction & 3.13 & 0.5103\\
& Error & 81.53 & \NA\\
\midrule
\multirow{4}{*}{Vigitel} & Separability & 96.00 & $5.03\times10^{-10}$\\
& Anchor & 0.08 & 0.6069\\
& Interaction & 0.33 & 0.3154\\
& Error & 3.59 & \NA\\
\bottomrule
\end{tabular}
\end{table}

\begin{figure}[htbp]
\centering
\includegraphics[width=.88\textwidth]{factorial_variation.pdf}
\caption{Variation allocation for the blocked raw-scale factorial models.}
\label{fig:variation}
\end{figure}

\begin{table}[htbp]
\centering
\caption{Residual assumption diagnostics for the blocked factorial models.}
\label{tab:assumptions}
\resizebox{\textwidth}{!}{%
\begin{tabular}{llrrrrp{2.5cm}}
\toprule
Dataset & Scale & Shapiro & B--F & Fligner & B--P & Assessment\\
\midrule
COVID & raw & .2446 & .8409 & .8394 & .4415 & no detected violation\\
COVID & log & .2031 & .8173 & .8050 & .6435 & no detected violation\\
Vigitel & raw & .8564 & .8413 & .7706 & .0186 & variance flag\\
Vigitel & log & .8625 & .7463 & .6834 & .0237 & variance flag\\
\bottomrule
\end{tabular}
}
\end{table}

Normality tests do not reject for either dataset. Vigitel's groupwise robust variance tests (Brown--Forsythe, Fligner--Killeen) do not reject equal variance, but Breusch--Pagan does, and the log transform does not resolve this conflict. More importantly, independence is not established, because the cross-validation training sets overlap and the four conditions within each fold share a candidate pool. The very small Vigitel residual should therefore be confirmed across independent generator and partition seeds before it is treated as a stable population effect (Section~\ref{sec:stat-limitations}).

\subsection{Summary of the empirical evidence}
\label{sec:results-summary}

The experiments form a coherent progression. Rebalancing and unconstrained augmentation did not make utility reliably high; domain filters sometimes helped but did not behave monotonically; diagnostic analyses implicated class overlap; \sepa{} then improved Vigitel \tstr{} performance while measured composite fidelity remained broadly stable; and the factorial study attributed that gain to separability rather than to anchors. The COVID-19 null result is equally informative: explicit separability pressure is not sufficient when the positive class is small, heterogeneous, or weakly represented.

The evidence gathered so far does not show that \sepa{} recovers clinically meaningful signal. An equally consistent explanation is that it deletes ambiguous observations and manufactures an easier classification problem (Section~\ref{sec:alternative}). Nor does the current factorial design establish generalisation beyond CTGAN, these two datasets, these two classifiers, or the recorded random seeds. We treat these as open questions rather than as incidental caveats, and return to them in Sections~\ref{sec:discussion} and~\ref{sec:limitations}.
